% =====================================================================
% preamble.tex — 宏包与全局设置
% 报告为英文，使用 pdflatex 编译（无需 XeLaTeX）
% =====================================================================

\usepackage[T1]{fontenc}
\usepackage[utf8]{inputenc}
\usepackage{lmodern}          % 更好的 Computer Modern 替代

% ---- 页面 ----
\usepackage[letterpaper, margin=1in]{geometry}

% ---- 数学 ----
\usepackage{amsmath}
\usepackage{amssymb}

% ---- 表格 ----
\usepackage{booktabs}        % \toprule \midrule \bottomrule
\usepackage{tabularx}
\usepackage{array}
\usepackage{multirow}
\usepackage{longtable}
\usepackage{seqsplit}        % 允许 \texttt 内断行，避免表格 Overfull

% ---- 图形 ----
\usepackage{graphicx}
\usepackage{tikz}            % 分类树（slides 与报告共用）
\usetikzlibrary{trees, positioning, arrows.meta, shapes.geometric, fit, backgrounds}

% ---- 参考文献（natbib + IEEEtran）----
\usepackage[numbers,sort&compress]{natbib}

% ---- 颜色（用于未核验标记与占位符）----
\usepackage{xcolor}

% ---- 超链接 ----
\usepackage[hidelinks, breaklinks=true]{hyperref}

% ---- 排版细节 ----
\usepackage{microtype}       % 更好的字距与断行
\usepackage{enumitem}        % 紧凑列表
\setlist{nosep}

% ---- 浮动体 ----
\usepackage[font=small, labelfont=bf, skip=4pt]{caption}
\usepackage{float}           % [H] 强制表格位置

% ---- 自定义命令 ----

% 论文 ID 引用，例如 \paperID{arxiv-2505.22954}
\newcommand{\paperID}[1]{\texttt{\footnotesize #1}}

% 可断行的行内代码，用于表格中的长 token
% 直接用 \texttt{failure-trajectory} 会撑爆列宽，必须用 \code{...}
\newcommand{\code}[1]{\texttt{\seqsplit{#1}}}

% 未核验标记
\newcommand{\unverified}{\textcolor{red}{\textbf{[UNVERIFIED]}}}

% 待填写占位
\newcommand{\todo}[1]{\textcolor{red}{\textit{[TODO: #1]}}}

% 空的表格单元格（灰色斜杠）
\newcommand{\na}{\textcolor{gray}{\textemdash}}

% 统一的表头样式
\newcommand{\thead}[1]{\textbf{\footnotesize #1}}

% =====================================================================
